\phantomsection
\subsection*{M1B. Semi-Sweet Mead}
\addcontentsline{toc}{subsection}{M1B. Semi-Sweet Mead}

\textbf{Impressão Geral}: Um hidromel com dulçor de mel perceptível e equilibrado. É percebido mais caráter de mel e corpo, quando comparado ao estilo M1A Dry Mead. Níveis de acidez e taninos sustentam o dulçor, mas o dulçor final não deve ser dominante. Apresenta equilíbrio, corpo, final e intensidade de sabor similar ao encontrado em uma ampla gama de vinhos semi-secos a doces/suaves.

\textbf{Aroma}: O aroma de mel deve ser de perceptível a moderado e pode ter um leve dulçor expresso como caráter de néctar de flores, reflexo de variedades de mel declaradas. O buquê de fermentação deve ser equilibrado. Hidroméis mais alcoólicos podem apresentar leves notas picantes oriundas do álcool. Os hidroméis carbonatados tendem a expressar mais aroma.

\textbf{Aparência}: Limpidez de alta a brilhante. Quanto maior o brilho, mais desejável. A presença de bolhas, espuma, colarinho ou efervescência deve corresponder ao nível de carbonatação declarado. Cor proveniente das variedades de méis declarados, geralmente de amarelo palha a dourado. Quanto mais vivas e brilhantes as cores, mais desejável.

\textbf{Sabor}: Caráter de mel de sutil a moderado, reflexo dos méis varietais declarados. Sabores semelhantes aos aromas de mel fermentado em Aroma. Níveis de dulçor residual são médio-baixos a médio-altos. Final com dulçor de baixo a moderado. A acidez deve ser equilibrada, macia, mas sem ser áspera ou agressiva. Taninos podem fazer um hidromel adocicado parecer mais seco. Características ásperas, desagradáveis ou excessivamente sulfurosas ou notas de levedura/autólise são indesejáveis.

\textbf{Sensação na Boca}: Corpo geralmente de médio-baixo a médio-alto. Hidroméis com maior força alcoólica podem apresentar corpo mais alto e maior aquecimento alcoólico. A sensação de corpo não deve ser acompanhada de dulçor residual maior que moderado. A carbonatação segue o nível declarado, variando de tranquilo a espumante.

\textbf{Ingredientes}: Mel de qualquer origem. Espera-se que os méis varietais apresentem características específicas associadas às variedades declaradas.

\textbf{Comentários}: Semi-Sweet Mead é caracterizado pelo dulçor equilibrado. Os níveis de dulçor podem ser variados, desde que estejam classificados entre hidroméis secos e suaves.

\textbf{Instruções para Inscrição}: Participantes devem obrigatoriamente declarar carbonatação e força alcoólica; dulçor é restrito a moderado nesse estilo. As variedades de mel podem ser declaradas.

\textbf{Exemplos Comerciais}: Colony Meadery Straight, No Chaser; Intermiel Verge d'Or; Rabbit's Foot Tupelonium; Redstone Traditional Mountain Honey Wine; Sky River Dry Mead
